% --- Archon physics formalization source begin ---
% archon:physics
% archon:covers PhyXMiniProblems/problem_phyx_mini_0782.lean
% archon:source-report reports/phyx_mini/problem_phyx_mini_0782.source.json

\chapter{Physics problem phyx\_mini\_0782}
\label{ch:PhyXMiniProblems_problem_phyx_mini_0782}

\paragraph{Problem source.}
\#\# Physical scenario

A thin, uniform, \textbackslash{}(3.80\textbackslash{}text\{ kg\}\textbackslash{}) bar, \textbackslash{}(80.0\textbackslash{}text\{ cm\}\textbackslash{}) long, has very small \textbackslash{}(2.50\textbackslash{}text\{ kg\}\textbackslash{}) balls glued on at either end as shown in figure, It is supported horizontally by a thin, horizontal, frictionless axle passing through its center and perpendicular to the bar. Suddenly the right-hand ball becomes detached and falls off, but the other ball remains glued to the bar.

\#\# Question

Find the angular velocity of the bar just as it swings through its vertical position.

\#\# Answer choices

- A: A: 4.60 rad/s

- B: B: 5.70 rad/s

- C: C: 6.20 rad/s

- D: D: 3.10 rad/s

\#\# Figure caption (auxiliary; use the image as primary evidence)

The image shows a horizontal bar with two spheres, each labeled "2.50 kg," positioned at each end. The bar has a labeled "Bar." In the middle, there's a black circle representing an "Axle (seen end-on)" that appears to be the pivot or rotation point. The structure implies symmetry, with the weights balanced on either side of the axle.

\#\# Dataset metadata

Rotational Motion; Physical Model Grounding Reasoning, Multi-Formula Reasoning

\paragraph{Recorded answer/context.}
B: B: 5.70 rad/s

\paragraph{Figure/image path.}
/root/proposal\_for\_physic/science-mango/phyx\_mini\_run/phyx\_data/test\_image/782.png

\paragraph{Formalization target.}
create a compiling Lean file with sorry bodies at `PhyXMiniProblems/problem\_phyx\_mini\_0782.lean`. The Lean declarations must preserve the physical quantities, dimensions or dimensional roles, figure labels, governing-law hypotheses, and final relation expressed by this problem.
Use Mathlib/Physlib names found through LeanExplore where available. If a physics API is missing, introduce faithful local abstractions rather than scalar placeholder aliases.

\begin{theorem}[Physics formalization target]
\label{thm:physics:phyx_mini_0782:target}
\lean{PhyXMiniProblems.ProblemPhyXMini0782.problem_phyx_mini_0782}
\uses{def:physics:phyx-mini-0782:phyxminiproblems-problemphyxmini0782-angularspeedinradianspersecond, def:physics:phyx-mini-0782:phyxminiproblems-problemphyxmini0782-centeraxlebarsetup, def:physics:phyx-mini-0782:phyxminiproblems-problemphyxmini0782-matchescenteredbarreleasescenario, def:physics:phyx-mini-0782:phyxminiproblems-problemphyxmini0782-matchesproblemreadouts, def:physics:phyx-mini-0782:phyxminiproblems-problemphyxmini0782-usesstandardterrestrialgravity, def:physics:phyx-mini-0782:phyxminiproblems-problemphyxmini0782-matchessuppliedfigure, def:physics:phyx-mini-0782:phyxminiproblems-problemphyxmini0782-haspositivephysicalparameters, def:physics:phyx-mini-0782:phyxminiproblems-problemphyxmini0782-satisfiescenteredbarinertialaws, def:physics:phyx-mini-0782:phyxminiproblems-problemphyxmini0782-satisfiesdetachedbarenergylaws, def:physics:phyx-mini-0782:phyxminiproblems-problemphyxmini0782-isuniqueclosestangularspeedchoice}
The assigned autoformalize agent should translate this physics problem into Lean declarations in the covered file, with theorem and lemma proof bodies written as `by sorry`.
\end{theorem}
\begin{proof}
This is an autoformalization task, not a proof task. Produce statements that can later be proved by the physics prover without weakening the physical model.
\end{proof}

% --- Archon declaration topology begin ---
\section{Declaration topology}
\begin{definition}[acceleration Dimension]
  \label{def:physics:phyx-mini-0782:phyxminiproblems-problemphyxmini0782-accelerationdimension}
  \lean{PhyXMiniProblems.ProblemPhyXMini0782.accelerationDimension}
  Gravitational acceleration has dimension length divided by time squared
\end{definition}
\begin{proof}
  This is the declared model definition of acceleration Dimension.
\end{proof}
\begin{definition}[moment Of Inertia Dimension]
  \label{def:physics:phyx-mini-0782:phyxminiproblems-problemphyxmini0782-momentofinertiadimension}
  \lean{PhyXMiniProblems.ProblemPhyXMini0782.momentOfInertiaDimension}
  Moment of inertia has dimension mass times length squared
\end{definition}
\begin{proof}
  This is the declared model definition of moment Of Inertia Dimension.
\end{proof}
\begin{definition}[Mass Quantity]
  \label{def:physics:phyx-mini-0782:phyxminiproblems-problemphyxmini0782-massquantity}
  \lean{PhyXMiniProblems.ProblemPhyXMini0782.MassQuantity}
  A nonnegative, unit-independent physical mass
\end{definition}
\begin{proof}
  This is the declared model definition of Mass Quantity.
\end{proof}
\begin{definition}[Length Quantity]
  \label{def:physics:phyx-mini-0782:phyxminiproblems-problemphyxmini0782-lengthquantity}
  \lean{PhyXMiniProblems.ProblemPhyXMini0782.LengthQuantity}
  A nonnegative, unit-independent physical length
\end{definition}
\begin{proof}
  This is the declared model definition of Length Quantity.
\end{proof}
\begin{definition}[Acceleration Quantity]
  \label{def:physics:phyx-mini-0782:phyxminiproblems-problemphyxmini0782-accelerationquantity}
  \lean{PhyXMiniProblems.ProblemPhyXMini0782.AccelerationQuantity}
  \uses{def:physics:phyx-mini-0782:phyxminiproblems-problemphyxmini0782-accelerationdimension}
  A nonnegative, unit-independent acceleration magnitude
\end{definition}
\begin{proof}
  This is the declared model definition of Acceleration Quantity.
\end{proof}
\begin{definition}[Angular Speed Quantity]
  \label{def:physics:phyx-mini-0782:phyxminiproblems-problemphyxmini0782-angularspeedquantity}
  \lean{PhyXMiniProblems.ProblemPhyXMini0782.AngularSpeedQuantity}
  A nonnegative angular-speed magnitude; radians are dimensionless
\end{definition}
\begin{proof}
  This is the declared model definition of Angular Speed Quantity.
\end{proof}
\begin{definition}[Moment Of Inertia Quantity]
  \label{def:physics:phyx-mini-0782:phyxminiproblems-problemphyxmini0782-momentofinertiaquantity}
  \lean{PhyXMiniProblems.ProblemPhyXMini0782.MomentOfInertiaQuantity}
  \uses{def:physics:phyx-mini-0782:phyxminiproblems-problemphyxmini0782-momentofinertiadimension}
  A nonnegative moment of inertia about the fixed axle
\end{definition}
\begin{proof}
  This is the declared model definition of Moment Of Inertia Quantity.
\end{proof}
\begin{definition}[Energy Quantity]
  \label{def:physics:phyx-mini-0782:phyxminiproblems-problemphyxmini0782-energyquantity}
  \lean{PhyXMiniProblems.ProblemPhyXMini0782.EnergyQuantity}
  Mechanical energy with Physlib's energy dimension
\end{definition}
\begin{proof}
  This is the declared model definition of Energy Quantity.
\end{proof}
\begin{definition}[mass Readout]
  \label{def:physics:phyx-mini-0782:phyxminiproblems-problemphyxmini0782-massreadout}
  \lean{PhyXMiniProblems.ProblemPhyXMini0782.massReadout}
  \uses{def:physics:phyx-mini-0782:phyxminiproblems-problemphyxmini0782-massquantity}
  Read a physical mass in a selected mass unit
\end{definition}
\begin{proof}
  This is the declared model definition of mass Readout.
\end{proof}
\begin{definition}[length Readout]
  \label{def:physics:phyx-mini-0782:phyxminiproblems-problemphyxmini0782-lengthreadout}
  \lean{PhyXMiniProblems.ProblemPhyXMini0782.lengthReadout}
  \uses{def:physics:phyx-mini-0782:phyxminiproblems-problemphyxmini0782-lengthquantity}
  Read a physical length in a selected length unit
\end{definition}
\begin{proof}
  This is the declared model definition of length Readout.
\end{proof}
\begin{definition}[acceleration Readout]
  \label{def:physics:phyx-mini-0782:phyxminiproblems-problemphyxmini0782-accelerationreadout}
  \lean{PhyXMiniProblems.ProblemPhyXMini0782.accelerationReadout}
  \uses{def:physics:phyx-mini-0782:phyxminiproblems-problemphyxmini0782-accelerationquantity}
  Read acceleration in coherent selected length and time units
\end{definition}
\begin{proof}
  This is the declared model definition of acceleration Readout.
\end{proof}
\begin{definition}[angular Speed Readout]
  \label{def:physics:phyx-mini-0782:phyxminiproblems-problemphyxmini0782-angularspeedreadout}
  \lean{PhyXMiniProblems.ProblemPhyXMini0782.angularSpeedReadout}
  \uses{def:physics:phyx-mini-0782:phyxminiproblems-problemphyxmini0782-angularspeedquantity}
  Read angular speed in inverse units of the selected time unit
\end{definition}
\begin{proof}
  This is the declared model definition of angular Speed Readout.
\end{proof}
\begin{definition}[moment Of Inertia Readout]
  \label{def:physics:phyx-mini-0782:phyxminiproblems-problemphyxmini0782-momentofinertiareadout}
  \lean{PhyXMiniProblems.ProblemPhyXMini0782.momentOfInertiaReadout}
  \uses{def:physics:phyx-mini-0782:phyxminiproblems-problemphyxmini0782-momentofinertiaquantity}
  Read moment of inertia in coherent selected mass and length units
\end{definition}
\begin{proof}
  This is the declared model definition of moment Of Inertia Readout.
\end{proof}
\begin{definition}[energy Readout]
  \label{def:physics:phyx-mini-0782:phyxminiproblems-problemphyxmini0782-energyreadout}
  \lean{PhyXMiniProblems.ProblemPhyXMini0782.energyReadout}
  \uses{def:physics:phyx-mini-0782:phyxminiproblems-problemphyxmini0782-energyquantity}
  Read energy in the coherent unit induced by the selected base units
\end{definition}
\begin{proof}
  This is the declared model definition of energy Readout.
\end{proof}
\begin{definition}[mass In Kilograms]
  \label{def:physics:phyx-mini-0782:phyxminiproblems-problemphyxmini0782-massinkilograms}
  \lean{PhyXMiniProblems.ProblemPhyXMini0782.massInKilograms}
  \uses{def:physics:phyx-mini-0782:phyxminiproblems-problemphyxmini0782-massquantity, def:physics:phyx-mini-0782:phyxminiproblems-problemphyxmini0782-massreadout}
  Kilogram readout of a physical mass
\end{definition}
\begin{proof}
  This is the declared model definition of mass In Kilograms.
\end{proof}
\begin{definition}[length In Meters]
  \label{def:physics:phyx-mini-0782:phyxminiproblems-problemphyxmini0782-lengthinmeters}
  \lean{PhyXMiniProblems.ProblemPhyXMini0782.lengthInMeters}
  \uses{def:physics:phyx-mini-0782:phyxminiproblems-problemphyxmini0782-lengthquantity, def:physics:phyx-mini-0782:phyxminiproblems-problemphyxmini0782-lengthreadout}
  Metre readout of a physical length
\end{definition}
\begin{proof}
  This is the declared model definition of length In Meters.
\end{proof}
\begin{definition}[acceleration In Meters Per Second Squared]
  \label{def:physics:phyx-mini-0782:phyxminiproblems-problemphyxmini0782-accelerationinmeterspersecondsquared}
  \lean{PhyXMiniProblems.ProblemPhyXMini0782.accelerationInMetersPerSecondSquared}
  \uses{def:physics:phyx-mini-0782:phyxminiproblems-problemphyxmini0782-accelerationquantity, def:physics:phyx-mini-0782:phyxminiproblems-problemphyxmini0782-accelerationreadout}
  Metre-per-second-squared readout of an acceleration magnitude
\end{definition}
\begin{proof}
  This is the declared model definition of acceleration In Meters Per Second Squared.
\end{proof}
\begin{definition}[moment Of Inertia In Kilogram Meters Squared]
  \label{def:physics:phyx-mini-0782:phyxminiproblems-problemphyxmini0782-momentofinertiainkilogrammeterssquared}
  \lean{PhyXMiniProblems.ProblemPhyXMini0782.momentOfInertiaInKilogramMetersSquared}
  \uses{def:physics:phyx-mini-0782:phyxminiproblems-problemphyxmini0782-momentofinertiaquantity, def:physics:phyx-mini-0782:phyxminiproblems-problemphyxmini0782-momentofinertiareadout}
  Kilogram-metre-squared readout of a moment of inertia
\end{definition}
\begin{proof}
  This is the declared model definition of moment Of Inertia In Kilogram Meters Squared.
\end{proof}
\begin{definition}[angular Speed In Radians Per Second]
  \label{def:physics:phyx-mini-0782:phyxminiproblems-problemphyxmini0782-angularspeedinradianspersecond}
  \lean{PhyXMiniProblems.ProblemPhyXMini0782.angularSpeedInRadiansPerSecond}
  \uses{def:physics:phyx-mini-0782:phyxminiproblems-problemphyxmini0782-angularspeedquantity, def:physics:phyx-mini-0782:phyxminiproblems-problemphyxmini0782-angularspeedreadout}
  Radian-per-second readout of angular-speed magnitude
\end{definition}
\begin{proof}
  This is the declared model definition of angular Speed In Radians Per Second.
\end{proof}
\begin{definition}[Bar End]
  \label{def:physics:phyx-mini-0782:phyxminiproblems-problemphyxmini0782-barend}
  \lean{PhyXMiniProblems.ProblemPhyXMini0782.BarEnd}
  The two ends of the horizontal bar in the primary figure
\end{definition}
\begin{proof}
  This is the declared model definition of Bar End.
\end{proof}
\begin{definition}[Swing Stage]
  \label{def:physics:phyx-mini-0782:phyxminiproblems-problemphyxmini0782-swingstage}
  \lean{PhyXMiniProblems.ProblemPhyXMini0782.SwingStage}
  The two instants relevant to the energy balance
\end{definition}
\begin{proof}
  This is the declared model definition of Swing Stage.
\end{proof}
\begin{definition}[Bar Orientation]
  \label{def:physics:phyx-mini-0782:phyxminiproblems-problemphyxmini0782-barorientation}
  \lean{PhyXMiniProblems.ProblemPhyXMini0782.BarOrientation}
  Orientation of the bar in the vertical plane of its swing
\end{definition}
\begin{proof}
  This is the declared model definition of Bar Orientation.
\end{proof}
\begin{definition}[Axle Location]
  \label{def:physics:phyx-mini-0782:phyxminiproblems-problemphyxmini0782-axlelocation}
  \lean{PhyXMiniProblems.ProblemPhyXMini0782.AxleLocation}
  Location of the axle along the bar
\end{definition}
\begin{proof}
  This is the declared model definition of Axle Location.
\end{proof}
\begin{definition}[Axle Friction Model]
  \label{def:physics:phyx-mini-0782:phyxminiproblems-problemphyxmini0782-axlefrictionmodel}
  \lean{PhyXMiniProblems.ProblemPhyXMini0782.AxleFrictionModel}
  Idealized friction model for the axle
\end{definition}
\begin{proof}
  This is the declared model definition of Axle Friction Model.
\end{proof}
\begin{definition}[Bar Mass Model]
  \label{def:physics:phyx-mini-0782:phyxminiproblems-problemphyxmini0782-barmassmodel}
  \lean{PhyXMiniProblems.ProblemPhyXMini0782.BarMassModel}
  Mass-distribution model for the bar
\end{definition}
\begin{proof}
  This is the declared model definition of Bar Mass Model.
\end{proof}
\begin{definition}[Endpoint Ball Model]
  \label{def:physics:phyx-mini-0782:phyxminiproblems-problemphyxmini0782-endpointballmodel}
  \lean{PhyXMiniProblems.ProblemPhyXMini0782.EndpointBallModel}
  Geometric idealization for either very small endpoint ball
\end{definition}
\begin{proof}
  This is the declared model definition of Endpoint Ball Model.
\end{proof}
\begin{definition}[Post Release Attachment]
  \label{def:physics:phyx-mini-0782:phyxminiproblems-problemphyxmini0782-postreleaseattachment}
  \lean{PhyXMiniProblems.ProblemPhyXMini0782.PostReleaseAttachment}
  Attachment state of an endpoint ball immediately after release
\end{definition}
\begin{proof}
  This is the declared model definition of Post Release Attachment.
\end{proof}
\begin{definition}[Figure Object]
  \label{def:physics:phyx-mini-0782:phyxminiproblems-problemphyxmini0782-figureobject}
  \lean{PhyXMiniProblems.ProblemPhyXMini0782.FigureObject}
  Objects literally visible in the supplied bitmap
\end{definition}
\begin{proof}
  This is the declared model definition of Figure Object.
\end{proof}
\begin{definition}[Supplied Bar Axle Figure]
  \label{def:physics:phyx-mini-0782:phyxminiproblems-problemphyxmini0782-suppliedbaraxlefigure}
  \lean{PhyXMiniProblems.ProblemPhyXMini0782.SuppliedBarAxleFigure}
  \uses{def:physics:phyx-mini-0782:phyxminiproblems-problemphyxmini0782-barend, def:physics:phyx-mini-0782:phyxminiproblems-problemphyxmini0782-figureobject}
  Presentation-level evidence transcribed from the primary bitmap
\end{definition}
\begin{proof}
  This is the declared model definition of Supplied Bar Axle Figure.
\end{proof}
\begin{definition}[Center Axle Bar Setup]
  \label{def:physics:phyx-mini-0782:phyxminiproblems-problemphyxmini0782-centeraxlebarsetup}
  \lean{PhyXMiniProblems.ProblemPhyXMini0782.CenterAxleBarSetup}
  \uses{def:physics:phyx-mini-0782:phyxminiproblems-problemphyxmini0782-massquantity, def:physics:phyx-mini-0782:phyxminiproblems-problemphyxmini0782-lengthquantity, def:physics:phyx-mini-0782:phyxminiproblems-problemphyxmini0782-accelerationquantity, def:physics:phyx-mini-0782:phyxminiproblems-problemphyxmini0782-angularspeedquantity, def:physics:phyx-mini-0782:phyxminiproblems-problemphyxmini0782-momentofinertiaquantity, def:physics:phyx-mini-0782:phyxminiproblems-problemphyxmini0782-energyquantity, def:physics:phyx-mini-0782:phyxminiproblems-problemphyxmini0782-barend, def:physics:phyx-mini-0782:phyxminiproblems-problemphyxmini0782-swingstage, def:physics:phyx-mini-0782:phyxminiproblems-problemphyxmini0782-barorientation, def:physics:phyx-mini-0782:phyxminiproblems-problemphyxmini0782-axlelocation, def:physics:phyx-mini-0782:phyxminiproblems-problemphyxmini0782-axlefrictionmodel, def:physics:phyx-mini-0782:phyxminiproblems-problemphyxmini0782-barmassmodel, def:physics:phyx-mini-0782:phyxminiproblems-problemphyxmini0782-endpointballmodel, def:physics:phyx-mini-0782:phyxminiproblems-problemphyxmini0782-postreleaseattachment, def:physics:phyx-mini-0782:phyxminiproblems-problemphyxmini0782-suppliedbaraxlefigure}
  The independent physical data for the bar-plus-one-ball system after the right ball detaches. Energy and inertia fields are not defined from the answer; the governing-law predicates below relate them to the setup
\end{definition}
\begin{proof}
  This is the declared model definition of Center Axle Bar Setup.
\end{proof}
\begin{definition}[Matches Centered Bar Release Scenario]
  \label{def:physics:phyx-mini-0782:phyxminiproblems-problemphyxmini0782-matchescenteredbarreleasescenario}
  \lean{PhyXMiniProblems.ProblemPhyXMini0782.MatchesCenteredBarReleaseScenario}
  \uses{def:physics:phyx-mini-0782:phyxminiproblems-problemphyxmini0782-centeraxlebarsetup}
  Qualitative mechanics and geometry stated in the problem
\end{definition}
\begin{proof}
  This is the declared model definition of Matches Centered Bar Release Scenario.
\end{proof}
\begin{definition}[Matches Problem Readouts]
  \label{def:physics:phyx-mini-0782:phyxminiproblems-problemphyxmini0782-matchesproblemreadouts}
  \lean{PhyXMiniProblems.ProblemPhyXMini0782.MatchesProblemReadouts}
  \uses{def:physics:phyx-mini-0782:phyxminiproblems-problemphyxmini0782-massinkilograms, def:physics:phyx-mini-0782:phyxminiproblems-problemphyxmini0782-lengthinmeters, def:physics:phyx-mini-0782:phyxminiproblems-problemphyxmini0782-angularspeedinradianspersecond, def:physics:phyx-mini-0782:phyxminiproblems-problemphyxmini0782-centeraxlebarsetup}
  Numerical data stated in the text, including release from rest
\end{definition}
\begin{proof}
  This is the declared model definition of Matches Problem Readouts.
\end{proof}
\begin{definition}[Uses Standard Terrestrial Gravity]
  \label{def:physics:phyx-mini-0782:phyxminiproblems-problemphyxmini0782-usesstandardterrestrialgravity}
  \lean{PhyXMiniProblems.ProblemPhyXMini0782.UsesStandardTerrestrialGravity}
  \uses{def:physics:phyx-mini-0782:phyxminiproblems-problemphyxmini0782-accelerationinmeterspersecondsquared, def:physics:phyx-mini-0782:phyxminiproblems-problemphyxmini0782-centeraxlebarsetup}
  Standard terrestrial gravity used to evaluate the multiple-choice data
\end{definition}
\begin{proof}
  This is the declared model definition of Uses Standard Terrestrial Gravity.
\end{proof}
\begin{definition}[Matches Supplied Figure]
  \label{def:physics:phyx-mini-0782:phyxminiproblems-problemphyxmini0782-matchessuppliedfigure}
  \lean{PhyXMiniProblems.ProblemPhyXMini0782.MatchesSuppliedFigure}
  \uses{def:physics:phyx-mini-0782:phyxminiproblems-problemphyxmini0782-massinkilograms, def:physics:phyx-mini-0782:phyxminiproblems-problemphyxmini0782-centeraxlebarsetup}
  Literal information visible in the supplied raster
\end{definition}
\begin{proof}
  This is the declared model definition of Matches Supplied Figure.
\end{proof}
\begin{definition}[Has Positive Physical Parameters]
  \label{def:physics:phyx-mini-0782:phyxminiproblems-problemphyxmini0782-haspositivephysicalparameters}
  \lean{PhyXMiniProblems.ProblemPhyXMini0782.HasPositivePhysicalParameters}
  \uses{def:physics:phyx-mini-0782:phyxminiproblems-problemphyxmini0782-massreadout, def:physics:phyx-mini-0782:phyxminiproblems-problemphyxmini0782-lengthreadout, def:physics:phyx-mini-0782:phyxminiproblems-problemphyxmini0782-accelerationreadout, def:physics:phyx-mini-0782:phyxminiproblems-problemphyxmini0782-momentofinertiareadout, def:physics:phyx-mini-0782:phyxminiproblems-problemphyxmini0782-centeraxlebarsetup}
  Positivity conditions for the nondegenerate physical system
\end{definition}
\begin{proof}
  This is the declared model definition of Has Positive Physical Parameters.
\end{proof}
\begin{definition}[Satisfies Centered Bar Inertia Laws]
  \label{def:physics:phyx-mini-0782:phyxminiproblems-problemphyxmini0782-satisfiescenteredbarinertialaws}
  \lean{PhyXMiniProblems.ProblemPhyXMini0782.SatisfiesCenteredBarInertiaLaws}
  \uses{def:physics:phyx-mini-0782:phyxminiproblems-problemphyxmini0782-massreadout, def:physics:phyx-mini-0782:phyxminiproblems-problemphyxmini0782-lengthreadout, def:physics:phyx-mini-0782:phyxminiproblems-problemphyxmini0782-momentofinertiareadout, def:physics:phyx-mini-0782:phyxminiproblems-problemphyxmini0782-centeraxlebarsetup}
  For a thin uniform bar pivoted through its center, I\_bar = M L² / 12. The remaining very small ball is a point mass a distance L/2 from the axle, and the total inertia is the sum of the two
\end{definition}
\begin{proof}
  This is the declared model definition of Satisfies Centered Bar Inertia Laws.
\end{proof}
\begin{definition}[Satisfies Detached Bar Energy Laws]
  \label{def:physics:phyx-mini-0782:phyxminiproblems-problemphyxmini0782-satisfiesdetachedbarenergylaws}
  \lean{PhyXMiniProblems.ProblemPhyXMini0782.SatisfiesDetachedBarEnergyLaws}
  \uses{def:physics:phyx-mini-0782:phyxminiproblems-problemphyxmini0782-massreadout, def:physics:phyx-mini-0782:phyxminiproblems-problemphyxmini0782-lengthreadout, def:physics:phyx-mini-0782:phyxminiproblems-problemphyxmini0782-accelerationreadout, def:physics:phyx-mini-0782:phyxminiproblems-problemphyxmini0782-angularspeedreadout, def:physics:phyx-mini-0782:phyxminiproblems-problemphyxmini0782-momentofinertiareadout, def:physics:phyx-mini-0782:phyxminiproblems-problemphyxmini0782-energyreadout, def:physics:phyx-mini-0782:phyxminiproblems-problemphyxmini0782-centeraxlebarsetup}
  Rotational kinetic energy is I ω² / 2. Because the bar's center of mass is at the axle, only the remaining endpoint ball loses gravitational potential energy, by falling through the distance L/2. The frictionless axle implies conservation of mechanical energy between the two stages
\end{definition}
\begin{proof}
  This is the declared model definition of Satisfies Detached Bar Energy Laws.
\end{proof}
\begin{definition}[Angular Speed Answer Choice]
  \label{def:physics:phyx-mini-0782:phyxminiproblems-problemphyxmini0782-angularspeedanswerchoice}
  \lean{PhyXMiniProblems.ProblemPhyXMini0782.AngularSpeedAnswerChoice}
  Answer labels in the order shown in the exercise
\end{definition}
\begin{proof}
  This is the declared model definition of Angular Speed Answer Choice.
\end{proof}
\begin{definition}[displayed Angular Speed Radians Per Second]
  \label{def:physics:phyx-mini-0782:phyxminiproblems-problemphyxmini0782-displayedangularspeedradianspersecond}
  \lean{PhyXMiniProblems.ProblemPhyXMini0782.displayedAngularSpeedRadiansPerSecond}
  \uses{def:physics:phyx-mini-0782:phyxminiproblems-problemphyxmini0782-angularspeedanswerchoice}
  The four displayed angular-speed values, in radians per second
\end{definition}
\begin{proof}
  This is the declared model definition of displayed Angular Speed Radians Per Second.
\end{proof}
\begin{definition}[Is Closest Angular Speed Choice]
  \label{def:physics:phyx-mini-0782:phyxminiproblems-problemphyxmini0782-isclosestangularspeedchoice}
  \lean{PhyXMiniProblems.ProblemPhyXMini0782.IsClosestAngularSpeedChoice}
  \uses{def:physics:phyx-mini-0782:phyxminiproblems-problemphyxmini0782-angularspeedanswerchoice, def:physics:phyx-mini-0782:phyxminiproblems-problemphyxmini0782-displayedangularspeedradianspersecond}
  A displayed answer choice is at least as close as every other choice
\end{definition}
\begin{proof}
  This is the declared model definition of Is Closest Angular Speed Choice.
\end{proof}
\begin{definition}[Is Unique Closest Angular Speed Choice]
  \label{def:physics:phyx-mini-0782:phyxminiproblems-problemphyxmini0782-isuniqueclosestangularspeedchoice}
  \lean{PhyXMiniProblems.ProblemPhyXMini0782.IsUniqueClosestAngularSpeedChoice}
  \uses{def:physics:phyx-mini-0782:phyxminiproblems-problemphyxmini0782-angularspeedanswerchoice, def:physics:phyx-mini-0782:phyxminiproblems-problemphyxmini0782-isclosestangularspeedchoice}
  The selected displayed answer is the unique closest one
\end{definition}
\begin{proof}
  This is the declared model definition of Is Unique Closest Angular Speed Choice.
\end{proof}
\begin{lemma}[vertical angular speed squared]
  \label{lem:physics:phyx-mini-0782:phyxminiproblems-problemphyxmini0782-vertical-angular-speed-squared}
  \lean{PhyXMiniProblems.ProblemPhyXMini0782.vertical_angular_speed_squared}
  \uses{def:physics:phyx-mini-0782:phyxminiproblems-problemphyxmini0782-angularspeedinradianspersecond, def:physics:phyx-mini-0782:phyxminiproblems-problemphyxmini0782-centeraxlebarsetup, def:physics:phyx-mini-0782:phyxminiproblems-problemphyxmini0782-matchescenteredbarreleasescenario, def:physics:phyx-mini-0782:phyxminiproblems-problemphyxmini0782-matchesproblemreadouts, def:physics:phyx-mini-0782:phyxminiproblems-problemphyxmini0782-usesstandardterrestrialgravity, def:physics:phyx-mini-0782:phyxminiproblems-problemphyxmini0782-haspositivephysicalparameters, def:physics:phyx-mini-0782:phyxminiproblems-problemphyxmini0782-satisfiescenteredbarinertialaws, def:physics:phyx-mini-0782:phyxminiproblems-problemphyxmini0782-satisfiesdetachedbarenergylaws}
  The standard inertia and energy laws give ω² = (m g L) / (M L² / 12 + m (L/2)²) = 3675/113. This is a derived conclusion, not a premise
\end{lemma}
\begin{proof}
  The conclusion follows from the governing relations and figure readouts named in the statement of vertical angular speed squared.
\end{proof}
\begin{lemma}[vertical angular speed exact]
  \label{lem:physics:phyx-mini-0782:phyxminiproblems-problemphyxmini0782-vertical-angular-speed-exact}
  \lean{PhyXMiniProblems.ProblemPhyXMini0782.vertical_angular_speed_exact}
  \uses{def:physics:phyx-mini-0782:phyxminiproblems-problemphyxmini0782-angularspeedinradianspersecond, def:physics:phyx-mini-0782:phyxminiproblems-problemphyxmini0782-centeraxlebarsetup, def:physics:phyx-mini-0782:phyxminiproblems-problemphyxmini0782-matchescenteredbarreleasescenario, def:physics:phyx-mini-0782:phyxminiproblems-problemphyxmini0782-matchesproblemreadouts, def:physics:phyx-mini-0782:phyxminiproblems-problemphyxmini0782-usesstandardterrestrialgravity, def:physics:phyx-mini-0782:phyxminiproblems-problemphyxmini0782-haspositivephysicalparameters, def:physics:phyx-mini-0782:phyxminiproblems-problemphyxmini0782-satisfiescenteredbarinertialaws, def:physics:phyx-mini-0782:phyxminiproblems-problemphyxmini0782-satisfiesdetachedbarenergylaws}
  Nonnegativity of angular speed selects the positive square root
\end{lemma}
\begin{proof}
  The conclusion follows from the governing relations and figure readouts named in the statement of vertical angular speed exact.
\end{proof}
% --- Archon declaration topology end ---
% --- Archon physics formalization source end ---
